Banco &
Considera-se instituição financeira, para efeito desta lei, a pessoa jurídica de direito público ou privado, que tenha como atividade principal ou acessória, cumulativamente ou não, a captação, intermediação ou aplicação de recursos financeiros (Vetado) de terceiros, em moeda nacional ou estrangeira, ou a custódia, emissão, distribuição, negociação, intermediação ou administração de valores mobiliários. Cite LEI No 7.492, DE 16 DE JUNHO DE 1986. Na definição da lei é denominado instituição financeira, mas popularmente é conhecido como Banco. 
\\ \hline

Convenio &
Convênios, contratos de repasse e termos de parceria são acordos feitos entre União e entidades governamentais dos demais entes da Federação, ou organizações não-governamentais, para transferência de recursos financeiros a serem utilizados na execução de um objetivo comum. Portanto, as transferências voluntárias e legais serão viabilizadas a partir dos mencionados instrumentos e outros congêneres. Cite url https://portaldatransparencia.gov.br/entenda-a-gestao-publica/convenios-e-outros-acordos Instrumento que, na ausência de legislação específica, dispõe sobre a transferência de recursos financeiros provenientes do Orçamento Fiscal e da Seguridade Social da União para a execução de programas, projetos e atividades de interesse recíproco e em regime de mútua colaboração; Cite DECRETO Nº 11.531, DE 16 DE MAIO DE 2023 
\\ \hline

DadosPessoaisAnonimizados &
Representação de informações pessoais que foram submetidas a processo de anonimização. 
\\ \hline

DadosPessoaisAnonimizadosExpandido &
Representação de informações pessoais que foram submetidas a processo de anonimização, mas possui atributos opcionais que detalham os dados. 
\\ \hline

EntidadePublica &
É a unidade de atuação dotada de personalidade jurídica;; Cite LEI Nº 9.784 , DE 29 DE JANEIRO DE 1999. 
\\ \hline

EsferaOrcamentaria &
Classificação que tem por finalidade identificar se a despesa ou a receita estão inseridas nos orçamentos fiscal, da seguridade social ou de investimento das empresas estatais. Cite url https://www.congressonacional.leg.br/legislacao-e-publicacoes/glossario-orcamentario/-/orcamentario/termo/classificacao\_por\_esfera\_orcamentaria 
\\ \hline

FonteOrcamentaria &
A classificação por fontes ou destinações de recursos (FR) tem como objetivo agrupar receitas que possuam as mesmas normas de aplicação na despesa. Em regra, as fontes ou destinações de recursos reúnem recursos oriundos de determinados códigos da classificação por natureza da receita orçamentária, conforme regras previamente estabelecidas. Por meio do orçamento público, essas fontes ou destinações são associadas a determinadas despesas de forma a evidenciar os meios para atingir os objetivos públicos. Como mecanismo integrador entre a receita e a despesa, o código de FR exerce um duplo papel no processo orçamentário. Para a receita orçamentária, esse código tem a finalidade de indicar a destinação de recursos para a realização de determinadas despesas orçamentárias. Para a despesa orçamentária, identifica a origem dos recursos que estão sendo utilizados. Cite MCASP - 11º Edicao: Parte I: Seção 5.1 
\\ \hline

MicroRegiao &
Recorte geográfico formado por municípios e vinculado a uma região no modelo. 
\\ \hline

Municipio &
Unidade territorial municipal relacionada à região e à microrregião correspondentes. 
\\ \hline

OrgaoPublico &
É a unidade de atuação integrante da estrutura da Administração direta e da estrutura da Administração indireta; Cite LEI Nº 9.784 , DE 29 DE JANEIRO DE 1999. 
\\ \hline

OrgaoPublicoOrcamentario &
Os órgãos orçamentários, por sua vez, correspondem a agrupamentos de unidades orçamentárias. As dotações são consignadas às unidades orçamentárias, responsáveis pela realização das ações. Cite MCASP - 11º Edicao: Parte Il: Seções 4.2.1 
\\ \hline

PessoaFisica &
Pessoa natural identificada no modelo por CPF e nome, podendo também ter NIS informado. 
\\ \hline

PessoaJuridica &
Entidade com personalidade jurídica identificada no modelo por CNPJ e nome. 
\\ \hline

PoderDaFederacao &
Classificação que identifica o Poder da Federação. 
\\ \hline

Regiao &
Unidade territorial usada para agrupar municípios no contexto geográfico da ontologia. 
\\ \hline

UnidadeDeOrgaoPublico &
É o agrupamento de serviços subordinados ao mesmo órgão ou repartição a que serão consignadas dotações próprias. 
\\ \hline

UnidadeGestora &
44 Portaria MTP 1.467/2022, Art.2º - Definições: VI - unidade gestora: entidade ou órgão único, de natureza pública, de cada ente federativo, abrangendo todos os poderes, órgãos e entidades autárquicas e fundacionais, que tenha por finalidade a administração, o gerenciamento e a operacionalização do RPPS, incluindo a arrecada. Cite Portaria MTP nº 1.467, de 2 junho de 2022: Capítulo I - Definições 
\\ \hline

UnidadeOrcamentaria &
Constitui unidade orçamentária o agrupamento de serviços subordinados ao mesmo órgão ou repartição a que serão consignadas dotações próprias. Cite LEI No 4.320, DE 17 DE MARÇO DE 1964 
\\ \hline
